\subsection{A rank-one curve that defeats every augmentation size test}
\label{sec:rank-one-augmentation-obstruction}

\paragraph{Finding A27 (the augmentation size test can fail over every extension).}
Let $u\in\mathbf F_4$ satisfy $u^2+u+1=0$, and let
$C/\mathbf F_4$ be the smooth projective model of
\[
                    y^2+y=ux^3+x+u/x.
\tag{RO.1}
\]
This is a genus-two curve of $2$-rank one.
For every $m\ge1$, put $Q=4^m$, $G_m=J_C(\mathbf F_Q)$
and $N_m=\#C(\mathbf F_Q)$.
For every odd prime $r$ with $\operatorname{ord}_r(2)=r-1$,
one has
\[
                         |G_m/rG_m|<N_m.
\tag{RO.2}
\]
In particular, no prime and constant extension of this
curve can satisfy either sufficient numerical criterion
of A26.  This statement does not exclude rational
class-number witnesses that fail those sufficient
inequalities, or geometric theta witnesses.

\begin{proof}
\emph{1. Exact Frobenius data.}
The two Artin--Schreier poles of (RO.1) have odd orders
$3$ and $1$, so the genus and $2$-rank formulas give
$g=2$ and $f_C=1$ \cite[Lemma~2.6]{PriesZhu2012}.
The two ramified points are rational.
Direct finite-field enumeration gives
\[
                 N_1=6,\qquad N_2=28.
\tag{RO.3}
\]
The accompanying independent checker enumerates the
solutions in $\mathbf F_4$ and $\mathbf F_{16}$ directly,
including the two ramified points.  It embeds $u$ as
$z^2+z$ in $\mathbf F_{16}=\mathbf F_2[z]/(z^4+z+1)$.
The trace-zero method in the retained 72-model search
gives the same counts by a separate enumeration.

If $P_4(T)$ is the Frobenius characteristic polynomial,
(RO.3) gives $a_1=N_1-5=1$ and
$a_2=(a_1^2-(17-N_2))/2=6$.
The genus-two functional equation therefore gives
\[
 \begin{aligned}
 P_4(T)&=T^4+T^3+6T^2+4T+16\\
       &=(T^2-T+4)(T^2+2T+4)=A(T)B(T).
 \end{aligned}
\tag{RO.4}
\]
Only this exact polynomial identity is used below;
no elliptic isogeny classification is needed.

Let $t_m$ and $w_m$ be the sums of the $m$th powers of
the roots of $A$ and $B$, respectively.  They are integers,
with recurrences
\[
\begin{gathered}
 t_0=2,\quad t_1=1,\quad t_m=t_{m-1}-4t_{m-2},\\
 w_0=2,\quad w_1=-2,\quad w_m=-2w_{m-1}-4w_{m-2}.
\end{gathered}
\tag{RO.5}
\]
Put $e_m=Q+1-t_m$ and $h_m=Q+1-w_m$.
The Frobenius formulas \cite[Theorem~19.1(b)]{MilneAV} yield
\[
             |G_m|=e_mh_m,\qquad N_m=Q+1-t_m-w_m.
\tag{RO.6}
\]
All roots of $A$ and $B$ have complex absolute value $2$,
so
\[
 e_m\le Q+1+2\sqrt Q,
 \qquad N_m\ge L_m:=4^m+1-4\cdot2^m.
\tag{RO.7}
\]
Moreover $t_m\equiv1\pmod4$ for $m\ge1$ by (RO.5),
and hence
\[
                               4\mid e_m.
\tag{RO.8}
\]

\emph{2. The rational torsion rank for an eligible prime.}
Write $\kappa_m(r)=\dim_{\mathbf F_r}G_m[r]$.
Finite abelian-group structure gives
$|G_m/rG_m|=r^{\kappa_m(r)}$, and
$0\le\kappa_m(r)\le4$.
Let $F$ be the $4$-power Frobenius acting on $J_C[r]$.
Then $G_m[r]=\ker(F^m-1)$.

For $r\ge5$, the polynomials $A,B$ are coprime modulo $r$,
because their resultant is $36$: their difference is $3T$
and $A(0)=4$.
The characteristic polynomial (RO.4), Cayley--Hamilton and
Bezout therefore split the four-dimensional torsion space
as two $F$-stable, two-dimensional spaces
$W_A\oplus W_B$ on which $A(F)=0$ and $B(F)=0$,
respectively.  This uses no semisimplicity assumption.
Since
\[
                         (T-2)B(T)=T^3-8,
\]
one has $F^3=8$ on $W_B$.
If $F^m$ fixes a nonzero vector in $W_B$, it follows that
$2^{3m}\equiv1\pmod r$.
For an eligible prime this means
\[
                              r-1\mid3m.
\tag{RO.9}
\]
If $\kappa_m(r)\ge2$ and $W_B$ contributes no fixed
vector, all of $W_A$ is fixed by $F^m$.
Taking its determinant gives $4^m\equiv1\pmod r$,
so $r-1\mid2m$.
Consequently
\[
 r\ge5,\quad\kappa_m(r)\ge2
                  \quad\Longrightarrow\quad r\le3m+1.
\tag{RO.10}
\]

\emph{3. Every extension degree $m\ge10$.}
At $m=10$,
\[
       (3m+1)^4=923521<1044481=L_{10}.
\]
Also
$L_{m+1}-4L_m=8\cdot2^m-3>0$, while for $m\ge10$
\[
 \frac{(3(m+1)+1)^4}{(3m+1)^4}
          \le(34/31)^4<2.
\]
Induction gives
\[
                            (3m+1)^4<L_m\le N_m
                            \qquad(m\ge10).
\tag{RO.11}
\]
For $r=3$, $|G_m/3G_m|\le81<(3m+1)^4<N_m$.
For $r\ge5$ and $\kappa_m(r)\ge2$, (RO.10)--(RO.11)
prove (RO.2).  If $\kappa_m(r)=0$, then
$|G_m/rG_m|=1<N_m$.

It remains to consider $\kappa_m(r)=1$.
If $r\mid e_m$, (RO.8) and the fact that $r$ is odd give
\[
 r\le e_m/4\le\frac{Q+1+2\sqrt Q}{4}
                    <Q+1-4\sqrt Q\le N_m.
\tag{RO.12}
\]
The strict middle inequality holds for $\sqrt Q\ge8$:
after multiplying by four it is
$3((\sqrt Q)^2-6\sqrt Q+1)>0$.
If $r\nmid e_m$, then $r\mid h_m$ by (RO.6), and
$F^m$ has a fixed vector on $W_B$.
Equation (RO.9) and (RO.11) again give $r<N_m$.
Since $|G_m/rG_m|=r$ in this case, (RO.2) follows.

\emph{4. The nine remaining extension degrees.}
The exact recurrences (RO.5) and trial division give the
following complete table.  The middle column lists
\emph{all} eligible odd primes dividing $|G_m|$.
The last column is the maximum of
$r^{\min\{4,v_r(|G_m|)\}}$ over those primes, or $1$
when there are none.  It is an upper bound on
$|G_m/rG_m|$ for every eligible prime, including those
not dividing $|G_m|$.
\[
\begin{array}{c|r|l|r}
m&N_m&\text{eligible prime divisors}&\text{upper bound}\\
\hline
1&6&\varnothing&1\\
2&28&3&9\\
3&60&19&19\\
4&256&3,5,13&13\\
5&996&\varnothing&1\\
6&3976&3,19&81\\
7&16764&\varnothing&1\\
8&66016&3,5,13&13\\
9&260340&19,181&361.
\end{array}
\tag{RO.13}
\]
Every entry in the last column is strictly smaller than
$N_m$.  The independent checker retains the full
factorizations, prime certificates by trial division,
and multiplicative orders of $2$, as well as the
point enumerations and both induction anchors.
Thus (RO.2) holds for the remaining degrees and hence
for all $m\ge1$.
\end{proof}

\paragraph{Scope of the obstruction.}
The result refutes universal success of the numerical
augmentation criterion, even on $2$-rank-one curves.
The conclusion of A26 itself is a sufficient conditional
statement and remains valid.
No class-number calculation for every actual cyclic
cover is asserted here, and the numerical failure does
not imply absence of a rational witness or of a geometric
primitive-zero cover.
